\unit{Exposición 221 -- Los esquemas de Hilbert}
\sid{221}
\src{249}
\seminares{13.er año, 1960/61, n.º 221\\Mayo de 1961}
\paperhead{TÉCNICAS DE CONSTRUCCIÓN Y TEOREMAS DE EXISTENCIA EN GEOMETRÍA ALGEBRAICA\\
IV. LOS ESQUEMAS DE HILBERT}%
  {por Alexander GROTHENDIECK}

\subsection*{Introducción}
\addcontentsline{toc}{subsection}{Introducción}

Las técnicas expuestas en [2], I y II, eran en lo esencial independientes de
toda hipótesis de proyectividad sobre los esquemas considerados.
Desafortunadamente, en la actualidad no permiten resolver los problemas de
existencia planteados en la Exposición II de [2]. En la presente exposición y
en la siguiente resolveremos estos problemas bajo hipótesis de proyectividad.
Las técnicas empleadas son típicamente proyectivas y prácticamente no hacen
uso de los resultados de I y II de [2]. Aquí construiremos los «esquemas de
Hilbert», destinados a sustituir el uso de las coordenadas de Chow, como se
indicó en [2], II, n.º 2. En la exposición siguiente, la teoría del paso al
cociente en esquemas desarrollada en III, unida a la teoría de los esquemas de
Hilbert, nos permitirá, por ejemplo, construir los esquemas de Picard
(definidos en [2], II, n.º 3) bajo condiciones bastante generales.

En resumen, puede decirse que ahora se dispone de una técnica de construcciones
proyectivas bastante satisfactoria, salvo que todavía
falta\footnote{Véase el aditivo al final de la exposición.} un teorema de paso
al cociente por grupos tales como el grupo proyectivo, que actúan «sin puntos
fijos» (cf. [2], III, n.º 8). La situación parece incluso algo mejor en
geometría analítica (cuando uno se limita al estudio de espacios analíticos
«proyectivos» sobre un espacio analítico dado), pues para los espacios
analíticos desaparece la dificultad del paso al cociente por un grupo que
actúa convenientemente. Por otra parte, tanto en geometría algebraica como en
geometría analítica quedaría por desarrollar una técnica de construcción
válida sin hipótesis de proyectividad.

\fsection{1}{Familias acotadas de haces: propiedades de permanencia}

Sean $k$ un cuerpo y $X$ un $k$-preesquema, de tipo finito para simplificar.
Para toda extensión $K/k$ se obtiene, pues, un $K$-preesquema
$X_K=X\otimes_k K$. Si $\mathcal F$ es un haz coherente sobre $X_K$ y $K'$ es
una extensión de $K$, entonces $\mathcal F\otimes_K K'=\mathcal F_{K'}$ es un
haz cuasicoherente sobre $X_K\otimes_K K'=X_{K'}$. Dicho esto, si $K$ y $K'$
son dos extensiones cualesquiera de $k$, $\mathcal F$ es un haz cuasicoherente
sobre $X_K$ y $\mathcal F'$
\src{250}
es un haz cuasicoherente sobre $X_{K'}$, diremos que $\mathcal F$ y
$\mathcal F'$ son \emph{equivalentes} si existen $k$-homomorfismos de $K$ y
$K'$ en una misma extensión $K''$ de $k$ tales que $\mathcal F_{K''}$ y
$\mathcal F'_{K''}$ sean isomorfos sobre $X_{K''}$. Esta es una relación de
equivalencia, y nos interesaremos por las clases de equivalencia de haces para
esta relación y por los conjuntos de clases de equivalencia. Obsérvese que, si
$X$ es de tipo finito sobre $k$,\footnote{El texto impreso lleva aquí $X_0$;
solo $X$ está definido en este pasaje.} toda clase de haces coherentes puede
definirse por un haz coherente sobre algún $X_K$, donde $K$ es una extensión
de tipo finito de $k$. Por consiguiente, en la definición de las clases de
haces coherentes puede uno limitarse a las extensiones algebraicamente cerradas
de $k$; también puede limitarse a una extensión algebraicamente cerrada fija
$\Omega$ de $k$, de grado de trascendencia infinito. Dos haces coherentes
$\mathcal F$ y $\mathcal F'$ sobre $X_\Omega$ son entonces equivalentes si y
solo si existe un $k$-automorfismo $\sigma$ de $\Omega$ tal que
$\mathcal F\otimes_\Omega(\Omega,\sigma)$ sea isomorfo a $\mathcal F'$.
Se comprueba así que hay una correspondencia biyectiva entre las clases de
haces coherentes para una u otra definición.

Sean $E$ y $E'$ dos conjuntos de clases de haces coherentes sobre $X$.
Consideremos las clases de todos los haces de la forma
$\mathcal F\otimes\mathcal F'$, donde $\mathcal F$ y $\mathcal F'$ son haces
coherentes sobre un mismo $X_K$, la clase de $\mathcal F$ pertenece a $E$ y la
de $\mathcal F'$ pertenece a $E'$. Se obtiene así un conjunto de clases de
haces coherentes, que denotaremos por $E\otimes E'$. Se definen del mismo modo
los $\operatorname{Tor}_i(E,E')$, etc. En general, a toda función $U$ que
asocia a cada sucesión $\mathcal F_1,\ldots,\mathcal F_n$ de $n$ haces
coherentes sobre un mismo $X_K$ un conjunto
$U(\mathcal F_1,\ldots,\mathcal F_n)$ de haces coherentes sobre $X_K$, y que
posee la evidente propiedad de compatibilidad con los isomorfismos de haces y
las imágenes inversas por cambio de base, se le asocia una función, denotada
por el mismo símbolo $U$, que asigna a cada sucesión de $n$ conjuntos
$E_1,\ldots,E_n$ de clases de haces coherentes un conjunto
$U(E_1,\ldots,E_n)$ de clases de haces coherentes.

Nuestro propósito en este apartado es dar una definición de ciertos conjuntos
de clases de haces, llamados \emph{acotados}, y demostrar que las operaciones
$U$ más usuales, aplicadas a conjuntos acotados, vuelven a dar conjuntos
acotados.

Sea $X$ un preesquema de tipo finito sobre un $S$ noetheriano. Para todo
$s\in S$, la fibra $X_s$ es un preesquema de tipo finito sobre $k(s)$, y
consideraremos clases de haces coherentes sobre $X_s$ en el sentido anterior.
Esto da sentido a la expresión «clase de haces coherentes sobre una fibra de
$X/S$» y a las expresiones análogas. Del mismo modo, procediendo por separado
sobre cada fibra, pueden seguir considerándose operaciones como
$E\otimes E'$, etc., que hacen corresponder a sistemas de conjuntos de
\src{251}
clases de haces coherentes sobre las fibras de $X/S$ otro conjunto de clases
de haces coherentes sobre las fibras de $X/S$.

\result{DEFINICIÓN}{1.1} Sea $E$ un conjunto de clases de haces coherentes
sobre las fibras de $X/S$. Se dice que $E$ está \emph{acotado} si existen un
preesquema $S'$ de tipo finito sobre $S$ y un haz coherente $\mathcal F'$ sobre
$X'=X\times_S S'$ tales que $E$ esté contenido en el conjunto de clases de
haces sobre las fibras de $X/S$ definido por $\mathcal F'$.

Por definición, este último hace corresponder a un $s\in S$ las clases de los
haces $\mathcal F'\otimes_{S'}k(s')$, donde $s'$ recorre los puntos de $S'$
situados sobre $s$ (de modo que $k(s')$ es una extensión de $k(s)$ y
$X_s\otimes_{k(s)}k(s')$ se identifica con la fibra
$X'\otimes_{S'}k(s')=X'_{s'}$ de $X'$ en $s'$). Puede decirse que las familias
acotadas son aquellas contenidas en una \emph{familia algebraica} de haces
coherentes parametrizada por un $S'$ de tipo finito sobre $S$.

Una unión finita de familias acotadas es acotada (tómese el preesquema suma de
los preesquemas de parámetros $S_i$ que definen las familias algebraicas
mayorantes). Un cambio de base $T\longrightarrow S$ transforma una familia
acotada relativa a $X/S$ en una familia acotada relativa a $X_T/T$, y el
recíproco es cierto si $T\longrightarrow S$ es sobreyectivo (o, más
generalmente, si su imagen contiene los $s$ que intervienen efectivamente en la
familia dada $E$ para $X/S$). Esto reduce, en principio, la determinación de
las familias acotadas al caso en que $S$ es el espectro de un álgebra de tipo
finito sobre el anillo de los enteros $\mathbf Z$.

Si $E$ y $E'$ son familias acotadas de clases de haces relativas a $X/S$,
entonces $E\otimes E'$ también está acotada. En efecto, supongamos que $E$ y
$E'$ están mayoradas, respectivamente, por las familias algebraicas definidas
por $T\longrightarrow S$ y $\mathcal F$ sobre $X_T$, y por
$T'\longrightarrow S$ y $\mathcal F'$ sobre $X_{T'}$. Se ve entonces que
$E\otimes E'$ está mayorada por la familia algebraica definida por
$T''=T\times_S T'\longrightarrow S$ y el haz $\mathcal F''$ sobre
$X_{T''}$, producto tensorial de las imágenes inversas de $\mathcal F$ y
$\mathcal F'$ sobre $X_{T''}$. Este razonamiento solo es válido porque el
funtor $\mathcal F\otimes\mathcal F'$ en $\mathcal F,\mathcal F'$ es exacto a
derecha y, por tanto, conmuta con la extensión de la base (y, en particular,
con el paso a las fibras). No es aplicable directamente a operaciones locales
tales como $\operatorname{Tor}_i(E,E')$, $\operatorname{Hom}(E,E')$ y
$\operatorname{Ext}^i(E,E')$. No obstante, puede demostrarse que estas
operaciones siguen transformando conjuntos acotados en conjuntos acotados,
procediendo como para $E\otimes E'$ pero utilizando además resultados del tipo
siguiente (todos ellos contenidos en [3], IV, 6.11): una familia acotada $E$
está siempre mayorada por una familia algebraica definida por un haz coherente
$\mathcal F$ sobre un $X_T$ ($T$ de tipo finito sobre $S$) que es plano
relativamente a $T$. (Se «corta en trozos» el espacio de parámetros inicial.)
Tales
\src{252}
propiedades de planitud de haces adecuados aseguran, en efecto, que operaciones
como $\operatorname{Tor}_i(\mathcal F,\mathcal F')$ conmuten con cualquier
cambio de base. El mismo método se aplica a operaciones de naturaleza global:
imágenes directas e imágenes directas derivadas de haces coherentes por
morfismos propios, $\operatorname{Ext}$ globales relativos a morfismos propios
(cf. [5], III, \S~6), etc. Todas estas operaciones transforman familias
acotadas de haces en familias acotadas de haces. (N. B. -- Aquí los
preesquemas sobre los que se toman los diversos haces pueden cambiar por las
operaciones del tipo considerado.)

Los dos enunciados siguientes se demuestran esencialmente mediante la misma
técnica de planitud; para la descomposición primaria sobre las fibras de un
morfismo de tipo finito, véase en particular [5], IV.

\result{PROPOSICIÓN}{1.2} Sean $E$ y $E'$ conjuntos acotados de clases de
haces sobre las fibras de $X/S$, donde se supone que $X$ es propio sobre $S$.
Entonces:

\medskip
\noindent (i) La familia de núcleos, conúcleos e imágenes de homomorfismos
$\mathcal F\longrightarrow\mathcal F'$, donde la clase de $\mathcal F$
pertenece a $E$ y la de $\mathcal F'$ pertenece a $E'$, está acotada.

\smallskip
\noindent (ii) La familia de haces $\mathcal F''$ que son extensiones de un
$\mathcal F$ por un $\mathcal F'$, donde la clase de $\mathcal F$ pertenece a
$E$ y la de $\mathcal F'$ pertenece a $E'$, está acotada.

Mediante un cambio de base conveniente puede suponerse que $E$ y $E'$ están
definidos, respectivamente, por haces coherentes $\mathcal G$ y $\mathcal G'$
sobre un $X_T/T$, donde $T$ es de tipo finito sobre $S$. Además, puede
suponerse que se satisfacen ciertas hipótesis de planitud que implican que la
formación de los haces
$f_{T*}(\operatorname{Hom}_{X_T}(\mathcal G,\mathcal G'))$ y
$\operatorname{Ext}^{1}_{f_T}(\mathcal G,\mathcal G')$ conmuta con el cambio
de base para un morfismo arbitrario $T'\longrightarrow T$. Asimismo, puede
suponerse que estos haces coherentes sobre $T$ son localmente libres. Sean
$T_0$ y $T_1$ los fibrados vectoriales sobre $T$ cuyos haces de gérmenes de
secciones son, respectivamente, los haces precedentes. Se define entonces
canónicamente un homomorfismo
$\mathcal G_{T_0}\longrightarrow\mathcal G'_{T_0}$ de haces coherentes sobre
$X_{T_0}$ y una extensión
\[
0\longrightarrow\mathcal G'_{T_1}\longrightarrow\mathcal G''
\longrightarrow\mathcal G_{T_1}\longrightarrow0
\]
de haces coherentes sobre $X_{T_1}$ que posee una propiedad universal
evidente.\footnote{El texto impreso invierte los términos extremos de esta sucesión. Se trata de extensiones de $\mathcal G$ por $\mathcal G'$, como indica $\operatorname{Ext}^1(\mathcal G,\mathcal G')$.} El segundo haz define una familia algebraica que mayora la familia
considerada en (ii). Lo mismo vale para el núcleo, el conúcleo y la imagen del
homomorfismo
\src{253}
precedente, cuya consideración demuestra (i) (siempre que se suponga que el
conúcleo es plano relativamente a $T_0$, caso al que se reduce de nuevo
cortando $T_0$ en trozos\ldots).

\result{PROPOSICIÓN}{1.3} Sea $E$ una familia acotada de clases de haces sobre
las fibras de $X/S$. Entonces las clases de los haces estructurales de los
$(\operatorname{supp}\mathcal F)_{\mathrm{red}}$, donde $\mathcal F$ es un
haz coherente sobre un $X_K$, con $K$ algebraicamente cerrado, cuya clase
pertenece a $E$, forman una familia acotada.

Aquí $(\operatorname{supp}\mathcal F)_{\mathrm{red}}$ denota el soporte de
$\mathcal F$ provisto de la estructura reducida inducida; es decir, su haz
estructural es el cociente de $\mathcal O_{X_K}$ por el mayor haz de ideales
que define $\operatorname{supp}\mathcal F$. Puede demostrarse el resultado
análogo a (1.3) para los haces deducidos canónicamente de $\mathcal F$ por la
teoría de la descomposición primaria; por ejemplo, los
$\mathcal F/\mathcal F_i$, donde los $\mathcal F_i$ son los subhaces de
$\mathcal F$ primarios para las componentes del soporte de $\mathcal F$ y
mínimos con esta propiedad; o los $\mathcal O_{X_K}/\mathfrak p$, donde
$\mathfrak p$ es un haz de ideales primos asociado a $\mathcal F$; o los
$\mathcal O_{X_K}/\mathfrak q$, donde $\mathfrak q$ es un haz de ideales
primarios asociado a una componente del soporte de $\mathcal F$ (siendo el
cuerpo de referencia algebraicamente cerrado).

\fsection{2}{Familias acotadas y polinomio de Hilbert}

En lo que sigue suponemos que $X$ es proyectivo sobre $S$ y está provisto de
un haz muy amplio, denotado $\mathcal O_X(1)$. Para toda extensión $K$ de una
extensión residual $k(s)$ de un punto $s$ de $S$, se considera sobre $X_K$ el
haz correspondiente $\mathcal O_{X_K}(1)$, que sigue siendo muy amplio.

A todo haz coherente $\mathcal F$ sobre $X_K$ se le asocia la función
\[
P_{\mathcal F}(n)=\text{característica de Euler--Poincaré de $\mathcal F(n)$
sobre $X_K$},
\]
que es un polinomio en el entero $n$, llamado \emph{polinomio de Hilbert} de
$\mathcal F$. Para valores grandes de $n$, $P(n)$ no es sino la dimensión de
$H^0(X_K,\mathcal F(n))$ sobre $K$, puesto que los
$H^i(X_K,\mathcal F(n))$ se anulan para $i>0$ y $n$ grande.

Supongamos ahora que $\mathcal F$ es un haz coherente sobre $X$, plano sobre
$S$. Entonces los polinomios de Hilbert de los haces $\mathcal F_s$ inducidos
sobre fibras $X_s$ pertenecientes a una misma componente conexa de $S$ son
todos iguales [5], III, \S~7. De ello se sigue (sin hipótesis de planitud) que
el conjunto de polinomios de Hilbert de los haces $\mathcal F_s$, $s\in S$,
es finito para todo haz coherente $\mathcal F$ sobre $X$.
\src{254}
Recordemos además que, si $\mathcal F$ es un haz coherente sobre $X$, es
isomorfo a un haz cociente de un haz $\mathcal O_X(-n)^N$, para $n$ y $N$
suficientemente grandes. Por tanto, los haces $\mathcal F_s$ inducidos sobre
las fibras también son cocientes del haz $\mathcal O(-n)^N$ sobre la fibra.\footnote{El texto impreso omite aquí el exponente $N$; la restricción a la fibra conserva la potencia directa de rango $N$ del haz precedente.}

De estas dos observaciones se deduce la parte «solo si» del teorema siguiente:

\result{TEOREMA}{2.1} Sea $X$ proyectivo sobre un $S$ noetheriano y sea
$\mathcal O_X(1)$ muy amplio sobre $X$ relativamente a $S$. Sea $E$ un
conjunto de clases de haces sobre las fibras de $X/S$. Para que $E$ esté
acotado, es necesario y suficiente\footnote{La exposición impresa dice
solamente «es necesario»; la fe de erratas de 1962 de la Exposición 221 prescribe «es
necesario y suficiente».} que satisfaga las condiciones siguientes:

\medskip
\noindent (a) Existe un haz coherente $\mathcal L$ sobre $X$ (que puede
suponerse de la forma $\mathcal O_X(-n)^N$) tal que $E$ esté contenido en la
familia de clases de haces coherentes cocientes de haces de la forma
$\mathcal L_K$.

\smallskip
\noindent (b) Los polinomios de Hilbert $P_{\mathcal F}$ de los haces
$\mathcal F$ cuya clase pertenece a $E$ son elementos de un mismo conjunto
finito de polinomios.

Queda por demostrar la parte «si», que será un caso particular de un resultado
más preciso. Para todo Módulo coherente $\mathcal F$ sobre un preesquema de
tipo finito sobre un cuerpo $K$, y todo entero $r$, sea $N_r$ el sub-Módulo de
$\mathcal F$ cuyas secciones sobre un abierto son las secciones de $\mathcal F$
sobre dicho abierto cuyo soporte tiene dimensión $<r$. Se tiene, pues,
$N_r=\mathcal F$ para $r>\dim\operatorname{supp}\mathcal F$, mientras que
$N_r=0$ para $r\leq0$, y se obtiene una filtración creciente finita de
$\mathcal F$ cuyos factores $N_{r+1}/N_r$\footnote{El texto impreso invierte
el cociente, escribiendo $N_r/N_{r+1}$, lo que es incompatible con la
filtración creciente que acaba de definirse.} tienen como ciclos primos
asociados exactamente los ciclos primos asociados a $\mathcal F$ que tienen
dimensión $r$. Pondremos
\[
\mathcal F_{(r)}=\mathcal F/N_r.
\]
Entonces los ciclos primos asociados a $\mathcal F_{(r)}$ son exactamente los
ciclos primos asociados a $\mathcal F$ que tienen dimensión $\geq r$; en
particular, es igual a $\mathcal F$ si y solo si los ciclos primos asociados a
$\mathcal F$ tienen dimensión $\geq r$. Con esta notación:

\result{TEOREMA}{2.2} Bajo las hipótesis de (2.1), sea $s$ un entero y
supongamos que $E$ satisface la condición (a) y la siguiente forma debilitada
de (b):

\medskip
\noindent $(b_s)$ Los polinomios de Poincaré $P_{\mathcal F}$ de los haces
$\mathcal F$ cuya clase pertenece a $E$ tienen acotados sus coeficientes en
grados $\geq s-1$.

En estas condiciones, los haces $\mathcal F_{(s)}$ (cuando la clase de
$\mathcal F$ recorre $E$) forman una familia acotada. Además, los coeficientes
de grado $s-2$ de los $P_{\mathcal F}$ están minorados.
\src{255}
Así pues:

\result{COROLARIO}{2.3} Supongamos que los haces $\mathcal F$ cuya clase
pertenece a $E$ son tales que todos sus ciclos primos asociados tienen
dimensión $d$ con $s\leq d\leq r$. Entonces, en (2.1) (b), basta considerar los
coeficientes de los $P_{\mathcal F}$ en grados comprendidos entre $s-1$ y $r$.

El resto de este apartado está dedicado a esbozar la demostración de (2.2).
Los lemas clave son los dos lemas siguientes, el primero de los cuales es bien
conocido (y resume el contenido matemático útil de las coordenadas de Chow):

\result{LEMA 2.4 (CHOW)}{} Consideremos los haces estructurales de subesquemas
$Y$ de fibras $X_K$ ($K$ recorre las extensiones algebraicamente cerradas de
cuerpos residuales de $S$), donde $Y$ es reducido y todas sus componentes
tienen la misma dimensión $r$ (considerándose $\mathcal O_Y$ como un haz
cociente de $\mathcal O_X$). Si los grados de los $Y$ están acotados, los $Y$
forman una familia acotada.

Aquí, el grado $a$ de $Y$ puede definirse del modo más cómodo mediante el
coeficiente del término principal de
\[
P_{\mathcal O_Y}=a\frac{n^r}{r!}+\cdots.
\]

\result{LEMA}{2.5} Sea $\mathcal L$ un haz coherente sobre $X$, y sea $E$ un
conjunto de clases de haces cocientes $\mathcal F$ de haces $\mathcal L_K$
($K$ una extensión residual de $S$). Supongamos que las fibras de $X$ sobre
$S$ tienen dimensión $\leq r$, y escribamos
\[
P_{\mathcal F}(n)=a_{\mathcal F}\frac{n^r}{r!}
+b_{\mathcal F}\frac{n^{r-1}}{(r-1)!}
+\text{términos de grado $<r-1$}.
\]
Entonces el coeficiente $a_{\mathcal F}$ está acotado, mientras que
$b_{\mathcal F}$ está minorado. Si $b_{\mathcal F}$ está acotado, la familia
de los $\mathcal F_{(r)}$ está acotada.

\emph{Demostración.} -- Tras sustituir $S$ por una unión de subesquemas de
$S$ que recubren $S$, puede suponerse que existe un morfismo finito
$f:X\longrightarrow\mathbf P^r_S$ tal que $\mathcal O_X(1)$ es isomorfo a la
imagen inversa de $\mathcal O_{\mathbf P^r_S}(1)$. Por tanto, para todo haz
coherente $\mathcal F$ sobre $X$ se tiene
$P_{\mathcal F}=P_{f_*(\mathcal F)}$. Por otra parte, se comprueba fácilmente
(mediante la técnica del apartado anterior) que un conjunto de haces
$\mathcal F$ sobre $X$ está acotado si y solo si lo está el conjunto de los
$f_*(\mathcal F)$. Finalmente,
\[
f_*(\mathcal F)_{(r)}=f_*(\mathcal F_{(r)}).
\]
\src{256}
Esto nos reduce al caso $X=\mathbf P^r_S$. Además, puede suponerse que
$\mathcal L=\mathcal O_{\mathbf P^r_S}(k)^s$ para ciertos $k$ y $s$. El
coeficiente $a_{\mathcal F}$ satisface
\[
0\leq a_{\mathcal F}\leq s
\]
y, por tanto, está acotado. Es equivalente decir que el coeficiente de
$n^{r-1}$ en los $P_{\mathcal F}(n)$ está minorado (respectivamente, acotado),
o decir lo mismo de
$P_{\mathcal F}(n-k)=P_{\mathcal F(-k)}(n)$. Nos reducimos así al caso
\[
\mathcal L=\mathcal O_{\mathbf P^r_S}^{s}.
\]

Consideremos la sucesión exacta
\[
0\longrightarrow N_r\longrightarrow\mathcal F
\longrightarrow\mathcal F_{(r)}\longrightarrow0,
\]
de donde
\[
P_{\mathcal F}=P_{\mathcal F_{(r)}}+P_{N_r}.
\]
Como el coeficiente de $n^{r-1}$ en $P_{N_r}$ es positivo (pues
$\dim\operatorname{supp}N_r\leq r-1$), se tiene
\[
b_{\mathcal F_{(r)}}\leq b_{\mathcal F}.
\]
Así pues, para demostrar las afirmaciones del lema podemos sustituir
$\mathcal F$ por $\mathcal F_{(r)}$; es decir, podemos suponer que los haces
cocientes $\mathcal F$ de $\mathcal L$ considerados son libres de torsión.

Como $\mathbf P^r_K$ es normal, se sigue que $\mathcal F$ es localmente libre
de rango $a=a_{\mathcal F}$ sobre un abierto $U=\mathbf P^r_K-Y$, donde $Y$
tiene codimensión $\geq2$. Por tanto, $\bigwedge^a\mathcal F$ es un haz sobre
$\mathbf P^r_K$ cuya restricción a $U$ es invertible; luego (como
$\mathbf P^r_K$ es regular y $Y$ tiene codimensión $\geq2$) es isomorfo a la
restricción de un haz invertible sobre $\mathbf P^r_K$, determinado salvo
isomorfismo. Este último es de la forma $\mathcal O_{\mathbf P^r_K}(d)$ para
un entero $d$ bien determinado. Como $\bigwedge^a\mathcal F$ es un cociente
de $\bigwedge^a\mathcal O_{\mathbf P^r_K}^{s}\simeq
\mathcal O_{\mathbf P^r_K}^{N}$\footnote{El texto impreso sustituye aquí el
exponente $s$ por un $n$ no definido; $s$ es el rango del haz libre que se
proyecta sobre $\mathcal F$.}, donde $N=\binom{s}{a}$, posee $N$ secciones
canónicas, y estas definen, por tanto, secciones de
\src{257}
$\mathcal O_{\mathbf P^r_K}(d)$ sobre $U$ que son restricciones de secciones
$s_i$ ($1\leq i\leq N$) de $\mathcal O_{\mathbf P^r_K}(d)$ (pues
$\mathbf P^r_K$ es normal y $Y$ tiene codimensión $\geq2$). Estas secciones
generan $\mathcal O_{\mathbf P^r_K}(d)$ en los puntos de $U$, de modo que no
son todas nulas; esto implica $d\geq0$. Por otra parte, un cálculo sencillo da
\[
b_{\mathcal F}=a_{\mathcal F}(r+1)/2+d.
\]
En particular, esto demuestra que $b_{\mathcal F}>0$, por lo que
$b_{\mathcal F}$ está minorado. Está acotado si y solo si $d$ lo está;
demostremos que, en tal caso, $\mathcal F$ pertenece a una familia acotada.
Podemos suponer que $a_{\mathcal F}$ y $b_{\mathcal F}$ son fijos, digamos
iguales a $a$ y $b$, de modo que $d$ es fijo. El conocimiento de las $N$
secciones $s_i$ de $\mathcal O_{\mathbf P^r_K}(d)$, es decir, de un
homomorfismo
\[
s:\bigwedge^a\mathcal L_K\longrightarrow\mathcal O_{\mathbf P^r_K}(d),
\]
permite recuperar $\mathcal F$ como la coimagen del homomorfismo compuesto
correspondiente
\[
\mathcal L_K\longrightarrow
\operatorname{Hom}\!\left(\bigwedge^{a-1}\mathcal L_K,
\bigwedge^a\mathcal L_K\right)\longrightarrow
\operatorname{Hom}\!\left(\bigwedge^{a-1}\mathcal L_K,
\mathcal O_{\mathbf P^r_K}(d)\right),
\]
donde el primero es el homomorfismo canónico procedente del producto exterior
y el segundo se deduce de $s$. La conclusión se sigue ahora de (1.2) (i).

Combinando los dos lemas precedentes se obtiene:

\result{LEMA}{2.6} Bajo las hipótesis preliminares de (2.1), supongamos que,
para todo $\mathcal F$, se tiene
\[
P_{\mathcal F}(n)=a_{\mathcal F}\frac{n^r}{r!}
+b_{\mathcal F}\frac{n^{r-1}}{(r-1)!}
+\text{términos de grado $<r-1$},
\]
y que los coeficientes $a_{\mathcal F}$ están acotados. Entonces los
coeficientes $b_{\mathcal F}$ están minorados. Si los $b_{\mathcal F}$ están
acotados, los $\mathcal F_{(r)}$ forman una familia acotada.

Podemos suponer que los cuerpos base $K$ de los haces $\mathcal F$ son
algebraicamente cerrados. Dótese cada
$\operatorname{supp}\mathcal F_{(r)}$, unión de las componentes de
$\operatorname{supp}\mathcal F$ de dimensión $r$, de la estructura reducida
inducida. Entonces los grados de los $\operatorname{supp}\mathcal F_{(r)}$
están acotados superiormente por $a$, por lo que, en virtud de (2.4), los
$\operatorname{supp}\mathcal F_{(r)}$ forman un conjunto acotado. Además, para
cada componente de $\operatorname{supp}\mathcal F_{(r)}$, la longitud de
$\mathcal F_{(r)}$ a lo largo de dicha componente es $\leq a$. Por tanto, si
$\mathcal J_{\mathcal F}$ es el ideal que define
$\operatorname{supp}\mathcal F_{(r)}$, entonces $\mathcal F_{(r)}$
\src{258}
puede considerarse como un Módulo sobre el subesquema $Y_{\mathcal F}$ de $X$
definido por $\mathcal J_{\mathcal F}^{a}$. Como en el lema anterior, también
podemos reducirnos al caso $\mathcal F=\mathcal F_{(r)}$, de modo que
$\mathcal F$ proviene de un Módulo sobre $Y_{\mathcal F}$. Los
$Y_{\mathcal F}$ corresponden a una familia acotada de Módulos cocientes de
los $\mathcal O_{X_K}$ y, por consiguiente, provienen de un subesquema cerrado
$Y$ de un esquema $X\times_S T$. Podemos entonces aplicar (2.5) a $Y/T$ y
$\mathcal L\otimes_XY$, lo que da la conclusión.

Ya podemos demostrar (2.2) por inducción sobre la cota superior $r$ de
$\dim\operatorname{supp}\mathcal F$. El enunciado es trivial para $r<0$;
supongamos, pues, que $r\geq0$ y que el enunciado ha sido demostrado para
$r'<r$. Por (2.6), los $\mathcal F_{(r)}$ forman un conjunto acotado; por
tanto, en virtud de (1.2) (i), también lo forman los núcleos de los
homomorfismos $\mathcal L_K\longrightarrow\mathcal F_{(r)}$. Existe, pues, un
Módulo coherente $\mathcal L'$ sobre $X$ tal que los núcleos en cuestión y,
por consiguiente, también los
\[
N_r(\mathcal F)=\operatorname{Ker}(\mathcal F\longrightarrow\mathcal F_{(r)}),
\]
son cocientes de los Módulos $\mathcal L'_K$. Como los
$\mathcal F_{(r)}$ están acotados, los $P_{\mathcal F_{(r)}}$ permanecen
acotados, y la fórmula
\[
P_{\mathcal F}=P_{\mathcal F_{(r)}}+P_{N_r}
\]
demuestra que los $P_{N_r}$ satisfacen la misma condición $(b_s)$ que los
$P_{\mathcal F}$. Por tanto, los $N_r$ satisfacen las condiciones (a) y
$(b_s)$, y por la hipótesis de inducción los $(N_r)_{(s)}$ permanecen acotados.
Ahora bien, $\mathcal F_{(s)}$ es una extensión de $\mathcal F_{(r)}$ por
$(N_r)_{(s)}$; luego, por (1.2) (ii), los $\mathcal F_{(s)}$ permanecen
acotados. Para la última afirmación de (2.2), llevamos simultáneamente la inducción sobre el acotamiento y sobre la cota inferior del coeficiente siguiente. Los núcleos acotados de los morfismos de fuente fija hacia $\mathcal F_{(s)}$ proporcionan un mismo haz coherente que se aplica sobreyectivamente sobre cada $N_s$, pero (1.2) (i) por sí solo no hace acotada la familia de estos cocientes. Reaplicamos por inducción la afirmación sobre los coeficientes a estos cocientes; la aditividad del polinomio de Hilbert y el Lema 2.6 dan entonces la cota inferior requerida para el coeficiente siguiente. Esto termina la demostración.

\fsection{3}{Esquemas de Hilbert: definición y teorema de existencia}

Sea $X$ un preesquema sobre otro preesquema $S$, y sea $\mathcal F$ un Módulo
cuasicoherente sobre $X$. Denotemos por
\[
\operatorname{Quot}(\mathcal F/X/S)
\]
el conjunto de Módulos cuasicoherentes cocientes de $\mathcal F$ que son
planos sobre $S$. Sea ahora $S'\longrightarrow S$ un morfismo de cambio de
base, y pongamos $X'=X\times_S S'$ y
$\mathcal F'=\mathcal F\otimes_{\mathcal O_S}\mathcal O_{S'}$. Así, $X'$ es
un preesquema sobre $S'$ provisto de un Módulo
\src{259}
$\mathcal F'$, y podemos considerar
$\operatorname{Quot}(\mathcal F'/X'/S')$. Pongamos
\[
\underline{\operatorname{Quot}}_{\mathcal F/X/S}(S')
=\operatorname{Quot}(\mathcal F'/X'/S')
\qquad\text{(donde $X'=X\times_SS'$)}.
\]
Si ahora $S''\longrightarrow S'$ es un $S$-morfismo, entonces
$X''=X\times_SS''$ es isomorfo a $X'\times_{S'}S''$, y $\mathcal F''$ es
isomorfo a $\mathcal F'\otimes_{\mathcal O_{S'}}\mathcal O_{S''}$. Como el
funtor de imagen inversa
\[
\mathcal G'\longmapsto
\mathcal G'\otimes_{\mathcal O_{S'}}\mathcal O_{S''}
\]
de la categoría de Módulos cuasicoherentes sobre $X'$ a la categoría de
Módulos cuasicoherentes sobre $X''$ es exacto a derecha y transforma Módulos
planos sobre $S'$ en Módulos planos sobre $S''$, hay una aplicación natural
\[
\operatorname{Quot}(\mathcal F'/X'/S')\longrightarrow
\operatorname{Quot}(\mathcal F''/X''/S''),
\]
por lo que $\underline{\operatorname{Quot}}_{\mathcal F/X/S}(S')$ es un
funtor contravariante de $S'$ (preesquema sobre $S$), con valores en la
categoría de conjuntos. En adelante supondremos que $X$ es proyectivo sobre un
$S$ noetheriano y que $\mathcal F$ es coherente; para simplificar, nos
limitaremos a los $S'$ sobre $S$ que sean localmente noetherianos.

\result{TEOREMA}{3.1} Bajo estas hipótesis, el funtor contravariante
$\underline{\operatorname{Quot}}_{\mathcal F/X/S}$ sobre la categoría de los
$S$-preesquemas localmente noetherianos es representable por un $S$-preesquema
$\operatorname{Quot}_{\mathcal F/X/S}$ que es unión disjunta de una sucesión
de $S$-esquemas proyectivos (en particular,
$\operatorname{Quot}_{\mathcal F/X/S}$ es localmente de tipo finito sobre
$S$).

Obtendremos tal descomposición como sigue. Sea $\mathcal O_X(1)$ un haz
invertible sobre $X$, muy amplio relativamente a $S$. Para todo polinomio
$P(n)$ con coeficientes racionales, sea
$\operatorname{Quot}^P(\mathcal F/X/S)$ el subconjunto de
$\operatorname{Quot}(\mathcal F/X/S)$ formado por los cocientes coherentes
$\mathcal G$ de $\mathcal F$ que son planos sobre $S$ y cuyo polinomio de
Hilbert en todo $s\in S$ es igual a $P$. Pongamos
\[
\underline{\operatorname{Quot}}^P_{\mathcal F/X/S}(S')
=\operatorname{Quot}^P(\mathcal F'/X'/S').
\]
Esto da un subfuntor de
$\underline{\operatorname{Quot}}_{\mathcal F/X/S}$. La propiedad de
invariancia de los polinomios de Hilbert recordada en el Apartado 2 implica lo
siguiente: el funtor $\underline{\operatorname{Quot}}_{\mathcal F/X/S}$ es
representable si y solo si lo son los
$\underline{\operatorname{Quot}}^P_{\mathcal F/X/S}$; en tal caso, el
$S$-preesquema representante $\operatorname{Quot}_{\mathcal F/X/S}$ es
isomorfo a la unión disjunta de los
$\operatorname{Quot}^P_{\mathcal F/X/S}$ que representan los funtores
$\underline{\operatorname{Quot}}^P_{\mathcal F/X/S}$. Así, (3.1) se deduce
del teorema siguiente:
\src{260}
\result{TEOREMA}{3.2} Con las notaciones anteriores, el funtor
$\underline{\operatorname{Quot}}^P_{\mathcal F/X/S}$ está representado por un
$S$-preesquema proyectivo $\operatorname{Quot}^P_{\mathcal F/X/S}$.

El resto de este apartado está dedicado a la demostración de (3.2).

Sea $\nu$ un entero. Para todo $S'$ sobre $S$, denotemos por $A_\nu(S')$ el
conjunto de cocientes $\mathcal G=\mathcal F'/\mathcal K$ de
$\mathcal F'=\mathcal F\otimes_{\mathcal O_S}\mathcal O_{S'}$ que son
coherentes y planos sobre $S'$, y que satisfacen las condiciones siguientes:
\begin{align*}
\text{a.}\quad &R^if'_*(\mathcal G(n))=0 &&\text{para $i>0$ y $n\geq\nu$},\\
\text{b.}\quad &R^if'_*(\mathcal K(n))=0 &&\text{para $i>0$ y $n\geq\nu$},\\
\text{c.}\quad &f'_*(\mathcal K(\nu+k))
=\mathcal S'_k f'_*(\mathcal K(\nu)) &&\text{para $k\geq0$}.
\end{align*}
Para la última relación, se supone que $X$ está escrito como el espectro primo
homogéneo de un álgebra graduada cuasicoherente $\mathcal S_*$ sobre $S$,
graduada positivamente y generada por $\mathcal S_1$, y se pone
$\mathcal S'_* =\mathcal S_*\otimes_{\mathcal O_S}\mathcal O_{S'}$, de modo
que $X'$ es el espectro primo homogéneo de $\mathcal S'_*$. Para demostrar
(3.2), se reduce fácilmente al caso $X=\mathbf P^r_S$ (pues $S$ es una unión
de abiertos $U$ tales que $X|U$ es un subesquema cerrado de $\mathbf P^r_U$ y
$\mathcal O_X(1)$ está inducido por $\mathcal O_{\mathbf P^r_U}(1)$), y al
caso en que $\mathcal F$ es de la forma
$\mathcal O_{\mathbf P^r_S}^{N}$ y es, por tanto, plano sobre $S$. En la
discusión precedente, los haces $\mathcal K$ son entonces también planos sobre
$S'$. Las relaciones de Künneth ([5], III, \S~7) demuestran que las condiciones
(a) y (b) son estables por cambio de base e implican que, para $n\geq\nu$, la
formación de $f'_*(\mathcal G(n))$ y $f'_*(\mathcal K(n))$ conmuta con la
extensión de la base (loc. cit.); por tanto, (c) también es estable por
extensión de la base.

En otras palabras, $A_\nu(S')$ es un funtor contravariante de $S'$, más
precisamente un subfuntor de
$A(S')=\underline{\operatorname{Quot}}^P_{\mathcal F/X/S}(S')$. Al variar
$\nu$, se obtiene una sucesión creciente de subconjuntos $A_\nu(S')$ de
$A(S')$ cuya unión es $A(S')$, por un teorema bien conocido de Serre [5], III,
\S~2. Supongamos ahora que $\mathcal G$ es un cociente coherente de
$\mathcal F'$, plano sobre $S'$, y que $s$ es un punto de $S'$ tal que el
cambio de base $\operatorname{Spec}(k(s))\longrightarrow S'$ da un cociente
$\mathcal G_s$ de $\mathcal F_s$ que satisface las condiciones (a), (b) y (c),
es decir, que pertenece a $A_\nu(\operatorname{Spec}(k(s)))$. Existe entonces
un entorno abierto $U$ de $s$ sobre el cual se satisfacen las mismas
condiciones para $\mathcal G|f'^{-1}(U)$; es decir, este cociente pertenece a
$A_\nu(U)$. Para (a) y (b), esto resulta del «teorema de las funciones
holomorfas» [5], III, \S~7; la condición (c) resulta del lema de Nakayama y del
hecho de que, en cualquier caso, se sabe que
$f'_*(\mathcal K(n+k))=\mathcal S'_k f'_*(\mathcal K(n))$
\src{261}
para $n$ suficientemente grande y $k\geq0$ (cf. [5], III, \S~2). De estas
observaciones se deduce lo siguiente (compárese [4], IV): el funtor $A$ es
representable si y solo si los funtores $A_\nu$ son representables; en tal
caso, el $S$-preesquema $Q$ que representa $A$ es unión creciente de
subesquemas abiertos $Q_\nu$ que representan los $A_\nu$.

Sea
\[
M_* =\bigoplus_{n\geq0}f_*(\mathcal F(n))=\mathcal S_*^{N},
\]
de modo que
\[
M'_* =M_*\otimes_{\mathcal O_S}\mathcal O_{S'}
=\bigoplus_{n\geq0}f'_*(\mathcal F'(n))=\mathcal S'_*{}^{N}.
\]
De (a) se sigue que

\noindent (a') $f'_*(\mathcal G(n))$ es localmente libre de rango $P(n)$ para
$n\geq\nu$, [5], III, \S~7; y de (b), con $i=1$, que

\noindent (a'') $f'_*(\mathcal G(n))$ es un Módulo cociente de $M'_n$.

Además, el conocimiento de este Módulo cociente para $n=\nu$ determina, por
(c), los sub-Módulos $f'_*(\mathcal K(n))$ de $M'_n$ para $n\geq\nu$ y, por
tanto, determina $\mathcal K$ y luego $\mathcal G$. Se obtiene así una
aplicación inyectiva
\[
A_\nu(S')\longrightarrow\operatorname{Grass}_{P(\nu)}(M'_\nu)
\]
de $A_\nu(S')$ en el conjunto de Módulos cocientes localmente libres de
$M'_\nu$ de rango $P(\nu)$ y, por tanto, un homomorfismo funtorial
\[
i_\nu:A_\nu(S')\longrightarrow
\underline{\operatorname{Grass}}_{P(\nu)}(M_\nu)(S'),
\]
donde el funtor de la derecha está representado por el esquema de Grassmann
$\operatorname{Grass}_{P(\nu)}(M_\nu)$ (compárese [4], V), que es proyectivo
sobre $S$. Afirmo:

\result{LEMA}{3.3} El funtor $A_\nu(S')$ es representable, y el morfismo
$Q_\nu\longrightarrow\operatorname{Grass}_{P(\nu)}(M_\nu)$ que representa el
homomorfismo $i_\nu$ es una inmersión (lo que implica que $Q_\nu$ es
cuasiproyectivo sobre $S$).

Esta afirmación equivale a la siguiente (compárese [4], IV). Supongamos dado
un Módulo cociente $N$ de $M'_\nu$, localmente libre de rango $P(\nu)$. Existe
entonces un subpreesquema $Z$ de $S'$ tal que, para todo preesquema localmente
noetheriano $T'$
\src{262}
sobre $S'$, la imagen inversa de $N$ sobre $T'$ pertenece a
$\operatorname{Im}A_\nu(T')$ si y solo si $T'\longrightarrow S'$ se factoriza
a través del subpreesquema $Z$. Cambiando las notaciones, podemos suponer que
$S'=S$; es decir, se nos da un cociente $N_\nu$ de $M_\nu$ por un sub-Módulo
$R_\nu$. Este proviene de un elemento de $A_\nu(S)$ si y solo si satisface las
dos condiciones siguientes:

\medskip
\noindent (i) $M_{\nu+k}/\mathcal S_kR_\nu$ es localmente libre de rango
$P(\nu+k)$ para $k\geq0$.

\smallskip
\noindent (ii) El subhaz $\mathcal K$ de $\mathcal F$ definido por el
sub-Módulo graduado
\[
R_* =\sum_{k\geq0}\mathcal S_kR_\nu
\]
de $M_*$ ([5], II, \S~3), y el cociente
$\mathcal G=\mathcal F/\mathcal K$, satisfacen las condiciones (a) y (b)
anteriores.

Estas condiciones son manifiestamente necesarias. Recíprocamente, si se
satisfacen, el haz $\mathcal G$ definido en (ii), al ser isomorfo al haz
asociado al $\mathcal S_*$-Módulo graduado $N_*$ que es la suma de los
$M_{\nu+k}/\mathcal S_kR_\nu$, es plano sobre $S$ (pues sus fibras son
sumandos directos de localizaciones de $N_*$\footnote{El texto impreso lleva $N$; aquí se localiza el Módulo graduado $N_*$ que acaba de definirse.} en ideales primos homogéneos de
$\mathcal S_*$), y tiene polinomio de Hilbert $P$ por (i). En virtud de (ii),
se satisfacen (a') y (a''), de modo que, para $n\geq\nu$, el homomorfismo
natural $N_n\longrightarrow f_*(\mathcal G(n))$ es un homomorfismo
sobreyectivo entre Módulos localmente libres del mismo rango y, por tanto, un
isomorfismo. Así,
\[
f_*(\mathcal K(\nu+k))=\mathcal S_kR_\nu
\quad\text{para todo $k\geq0$}.
\]
Esto demuestra que $\mathcal G\in A_\nu(S)$ y que $R_\nu$ es el elemento de
$\operatorname{Grass}_{P(\nu)}(M_\nu)$ definido por $\mathcal G$.

El criterio (i)--(ii) también se aplica después de un cambio de base
$S'\longrightarrow S$. Demostraremos primero que la condición (i) después del
cambio de base $S'\longrightarrow S$ puede expresarse exigiendo que
$S'\longrightarrow S$ se factorice a través de cierto subpreesquema $Z$ de
$S$. Una vez hecho esto, sustituir $S$ por $Z$ nos reduce al caso en que (i) ya
se satisface sobre $S$; puesto que es estable por cambio de base, queda por
expresar la condición (ii). Sea $U$ el conjunto de los $s\in S$ tales que la
cohomología en dimensión $>0$ de los haces inducidos sobre la fibra $X_s$ por
$\mathcal G(n)$ y $\mathcal K(n)$ se anula para $n\geq\nu$. Ya hemos observado
que $U$ es abierto, y la condición (ii) se satisface tras un cambio de base
$S'\longrightarrow S$ si y solo si $S'\longrightarrow S$ se factoriza a
través de $U$. Esto demuestra el Lema 3.3. Queda por demostrar el lema
siguiente:

\result{LEMA}{3.4} Sea $S$ un preesquema localmente noetheriano provisto de un
álgebra cuasicoherente graduada positivamente $\mathcal S_*$, generada por
$\mathcal S_1$; sea $M_*$ un $\mathcal S_*$-Módulo graduado cuasicoherente de
tipo finito, $P$ un polinomio con coeficientes racionales y $\nu$ un entero.
Existe entonces un subpreesquema manifiestamente único $Z$ de $S$
\src{263}
con la propiedad siguiente: para todo preesquema $S'$ sobre $S$, el Módulo
$M_n\otimes_{\mathcal O_S}\mathcal O_{S'}$ es localmente libre de rango $P(n)$
para todo $n\geq\nu$ si y solo si $S'\longrightarrow S$ se factoriza a través
de $Z$.

Podemos suponer claramente que $S$ es afín y, por tanto, noetheriano. Entonces:

\result{LEMA}{3.5} Para todo entero $N\geq\nu$, sea $U_N$ el subconjunto
abierto de $S$ formado por los $s\in S$ tales que
\[
\operatorname{rank}_{k(s)}
M_{ns}\otimes_{\mathcal O_{S,s}}k(s)\leq P(n)
\quad\text{para todo $\nu\leq n\leq N$}.
\]
Entonces la sucesión decreciente de abiertos $U_N$ es estacionaria.

Se sabe [3], IV, que $S$ posee una partición finita en subesquemas reducidos
$S_i$ tales que cada $M\otimes_{\mathcal O_S}\mathcal O_{S_i}$ es plano sobre
$S_i$. Podemos, pues, suponer que $M$ es plano y, por tanto, que los $M_n$ son
planos. Finalmente, podemos suponer claramente que $S$ es conexo. Existen
entonces un entero $n_0$ y un polinomio $Q$ tales que
\[
\operatorname{rank}_{k(s)}
M_{ns}\otimes_{\mathcal O_{S,s}}k(s)=Q(n)
\quad\text{para $n\geq n_0$},
\]
[5], III, \S~7. Supongamos primero que $P<Q$\footnote{El texto impreso lleva
$P\neq Q$; la fe de erratas de 1962 prescribe $P<Q$, sustituye el caso «$P=Q$» por
el caso complementario $P(n)\geq Q(n)$ para $n$ grande, y sustituye las
desigualdades estrictas de (*) por desigualdades no estrictas.}, de modo que
$P(n)<Q(n)$ para $n$ grande. Entonces, evidentemente, $U_N=\varnothing$ para
$N$ grande, y la sucesión $U_N$ es estacionaria. En el caso complementario,
$P(n)\geq Q(n)$ para $n$ grande, de modo que $U_N=U_{n_0}$ para $N\geq n_0$,
y la sucesión $U_N$ vuelve a ser estacionaria.

En particular, el conjunto $U_\infty$ de los $s\in S$ tales que
\[
\tag{*}
\operatorname{rank}_{k(s)}
M_{ns}\otimes_{\mathcal O_{S,s}}k(s)\leq P(n)
\quad\text{para todo $n\geq\nu$}
\]
es abierto, pues es la intersección de los $U_N$. Para demostrar (3.4),
podemos, por tanto, sustituir $S$ por el abierto $U_\infty$\footnote{El texto impreso lleva $U$; el abierto considerado aquí es $U_\infty$, la intersección de los $U_N$.}, reduciéndonos al caso en
que la desigualdad (*) se satisface en todo $s\in S$. Además:

\result{LEMA}{3.6} Sea $M$ un Módulo sobre un preesquema localmente
noetheriano $S$, y sea $r$ un entero. Existe entonces un subpreesquema
manifiestamente único $Z$ de $S$ con la propiedad siguiente: para todo $S'$
sobre $S$, el Módulo $M\otimes_{\mathcal O_S}\mathcal O_{S'}$ es localmente
libre de rango $r$ si y solo si $S'\longrightarrow S$ se factoriza a través de
$Z$. Si
\[
\operatorname{rank}_{k(s)}
M_s\otimes_{\mathcal O_{S,s}}k(s)\leq r
\quad\text{para todo $s$},
\]
entonces $Z$ es un subpreesquema cerrado de $S$ (suponiéndose $M$ coherente).
\src{264}
En efecto, el razonamiento anterior nos reduce al caso en que la desigualdad
(*) se satisface para todo $s\in S$ (sustituyendo $S$, si es necesario, por el
subconjunto abierto de los puntos donde se satisface). Podemos entonces suponer
que $M$ interviene en una sucesión exacta
\[
\mathcal O_S^q\longrightarrow\mathcal O_S^r\longrightarrow M\longrightarrow0.
\]
La condición considerada para $S'$ sobre $S$ significa también que, en la
sucesión exacta correspondiente
$\mathcal O_{S'}^q\longrightarrow\mathcal O_{S'}^r
\longrightarrow M'\longrightarrow0$, la segunda flecha es un isomorfismo; es
decir, que la primera es nula. Por consiguiente, el subpreesquema cerrado $Z$
de $S$ definido por el ideal generado por las entradas de la matriz que define
$\mathcal O_S^q\longrightarrow\mathcal O_S^r$ posee la propiedad requerida.

Volvamos ahora a la demostración de (3.4). Sea $Z_n$ el subesquema cerrado de
$S$ asociado por (3.6) al Módulo $M_n$ y al entero $r=P(n)$, y sea $Z'_N$ el
ínfimo de los $Z_n$ para $\nu\leq n\leq N$. Los $Z'_N$ forman una sucesión
decreciente de subesquemas cerrados de $S$\footnote{El texto impreso lleva $Z$; todos los $Z_n$ y sus ínfimos $Z'_N$ son subesquemas cerrados de $S$.} y, por tanto, es necesariamente
estacionaria. Sea $Z$ su valor constante para $N$ grande. Este es el $Z$
buscado en (3.4). Esto termina la demostración de (3.4) y, por consiguiente,
también la de (3.3).

Hemos demostrado, pues, que
$\underline{\operatorname{Quot}}^P_{\mathcal F/X/S}$ está representado por un
$S$-preesquema $Q$ que es unión creciente de subesquemas abiertos $Q_\nu$
cuasiproyectivos sobre $S$. Para ir más lejos, invoquemos el Teorema 2.1, del
cual se deduce fácilmente que $Q$ es cuasicompacto (pues es la imagen del
preesquema $S'$ de tipo finito sobre $S$ que parametriza la familia de haces
cocientes de los $\mathcal F_K$ con polinomio de Hilbert $P$). Así, $Q$ es
igual a uno de los $Q_\nu$ y, por tanto, es cuasiproyectivo sobre $S$. Para
demostrar que es proyectivo sobre $S$, queda por demostrar que es propio sobre
$S$. Para ello basta invocar el criterio valuativo de la propiedad de ser
propio en la forma
[5], II, 7.3.8. Solo hay que verificar:

\result{LEMA}{3.7} Sea $S$ el espectro de un anillo de valoración discreta,
sea $s$ su punto genérico, sea $X$ un preesquema sobre $S$, sea $\mathcal F$
un Módulo cuasicoherente sobre $X$, y sea $\mathcal G_s$ un Módulo
cuasicoherente cociente de
$\mathcal F_s=\mathcal F\otimes_{\mathcal O_S}k(s)$ sobre $X_s$. Existe
entonces un único Módulo cuasicoherente cociente $\mathcal G$ de $\mathcal F$,
plano sobre $S$, cuya restricción a $X_s$ es $\mathcal G_s$.
\src{265}
En efecto, si $\mathcal G_s=\mathcal F_s/\mathcal K_s$, basta tomar el mayor
subhaz $\mathcal K$ de $\mathcal F$ que induce $\mathcal K_s$, [5], I, 9.4.2,
y poner $\mathcal G=\mathcal F/\mathcal K$. Se comprueba fácilmente que este
haz posee las propiedades requeridas.

El Teorema 3.2 y, por consiguiente, (3.1), quedan ahora completamente
demostrados.

La demostración establece al mismo tiempo:

\result{PROPOSICIÓN}{3.8} Bajo las hipótesis de (3.2), sean
$Q=\operatorname{Quot}^P_{\mathcal F/X/S}$, $X_Q=X\times_SQ$ y
$\mathcal F_Q=\mathcal F\otimes_{\mathcal O_S}\mathcal O_Q$; y sea
$\mathcal G$ el cociente coherente de $\mathcal F_Q$, plano sobre $Q$, con
polinomio de Hilbert relativo $P$, tal que $(Q,\mathcal G)$ representa el
funtor $\underline{\operatorname{Quot}}^P_{\mathcal F/X/S}$. Existe un entero
$\nu$ tal que, para $n\geq\nu$, $(f_Q)_*(\mathcal G(n))$ es un Módulo sobre
$Q$, localmente libre de rango $P(n)$ y muy amplio relativamente a $S$; es
decir, define una inmersión de $Q$ en un esquema de Grassmann
$\operatorname{Grass}_{P(n)}(M)$ sobre $S$. En particular, para $n\geq\nu$,
el haz
\[
\bigwedge^{P(n)}(f_Q)_*(\mathcal G(n))
\]
sobre $Q$ es invertible y muy amplio relativamente a $S$.

Como en (3.2), podemos reducirnos al caso en que $\mathcal F$ es plano sobre
$S$; entonces basta elegir un entero $\nu$ tal que $A_\nu=A$ (con las
notaciones anteriores).

La aplicación más importante de (3.2) es el caso
$\mathcal F=\mathcal O_X$. Escribimos entonces
\[
\operatorname{Quot}_{\mathcal O_X/X/S}=\operatorname{Hilb}_{X/S},
\qquad
\operatorname{Quot}^P_{\mathcal O_X/X/S}=\operatorname{Hilb}^P_{X/S},
\]
y obtenemos una descomposición
\[
\operatorname{Hilb}_{X/S}=\coprod_P\operatorname{Hilb}^P_{X/S}.
\]
Por definición, $\operatorname{Hilb}_{X/S}$ representa el funtor
$\underline{\operatorname{Hilb}}_{X/S}(S')$, conjunto de los subesquemas
cerrados de $X'=X\times_SS'$ que son planos sobre $S'$; y
$\operatorname{Hilb}^P_{X/S}$ representa el subfuntor correspondiente de
subesquemas cerrados con polinomio de Hilbert dado $P$. Estos preesquemas se
llaman también, respectivamente, \emph{preesquema de Hilbert} de $X$ sobre
$S$ y \emph{preesquema de Hilbert de índice $P$}. La terminología se justifica
por el papel desempeñado en la teoría por los polinomios de Hilbert. Su
diferencia de naturaleza respecto de las variedades de Chow clásicas (que
parametrizan ciclos y no variedades) es del mismo orden que la existente entre
\src{266}
el anillo de Chow de las clases de ciclos sobre una variedad y el anillo de
clases de haces sobre la variedad (tal como aparece en el teorema de
Riemann--Roch [1]). En efecto, cuando $X=\mathbf P^r_S$ y $S$ es el espectro
de un cuerpo, conocer el polinomio de Hilbert de un Módulo coherente
$\mathcal F$ sobre $X$ equivale también a conocer las clases de Chern de
$\mathcal F$, o la clase de $\mathcal F$ en el anillo de clases de haces
coherentes sobre $X$.

\result{OBSERVACIONES}{3.9} Obsérvese asimismo que la construcción de
$\operatorname{Quot}_{\mathcal F/X/S}$ y
$\operatorname{Quot}^P_{\mathcal F/X/S}$ se redujo al caso
$X=\mathbf P^r_S$, $\mathcal F=\mathcal O_X^N$, con $\mathcal O_X(1)$ el haz
muy amplio usual; más precisamente, los
$\operatorname{Quot}_{\mathcal F/X/S}$ generales aparecen como subesquemas
cerrados de estos. Como la formación de los
$\operatorname{Quot}^P_{\mathcal F/X/S}$ es manifiestamente compatible con
los cambios de base $S'\longrightarrow S$, nos reducimos además a
$S=\operatorname{Spec}(\mathbf Z)$. En definitiva, pues, se reduce al estudio
de los esquemas proyectivos siguientes sobre $\mathbf Z$:
\[
\mathfrak Q^P_{r,N}=
\operatorname{Quot}^{P}_{(\mathcal O_{\mathbf P^r_{\mathbf Z}})^N/
\mathbf P^r_{\mathbf Z}/\operatorname{Spec}(\mathbf Z)},
\]
y, más particularmente, de los esquemas de Hilbert absolutos
\[
\operatorname{Hilb}^P_r=\mathfrak Q^P_{r,1}.
\]
Sería muy interesante un estudio más detallado de estos esquemas, comenzando
por la determinación de sus componentes conexas (¿son conexos?),\footnote{La
fe de erratas de 1962 señala que Hartshorne demostró la conexidad de los
$\operatorname{Hilb}^P_r$ sobre un cuerpo algebraicamente cerrado y determinó
los pares $(r,P)$ para los cuales
$\operatorname{Hilb}^P_r\neq\varnothing$.} y de sus componentes irreducibles
(por Serre [7], puede haber componentes irreducibles situadas enteramente sobre
un primo no nulo $p$). Recordemos la pregunta de Weil de si las componentes
irreducibles de las fibras de $\operatorname{Hilb}^P_r$ sobre
$s\in\operatorname{Spec}(\mathbf Z)$ corresponden a extensiones «regulares»
del cuerpo primo, es decir, si son «relativamente conexas». Estas cuestiones
quizá sean más abordables para los esquemas de Hilbert que para las
«variedades de Chow».

\fsection{4}{Variantes}

\noindent a. Bajo las hipótesis de (3.1), sea $U$ un subconjunto abierto de
$X$, y sea $A'$ el subfuntor de
$A=\underline{\operatorname{Quot}}_{\mathcal F/X/S}$ tal que $A'(S')$ es el
conjunto de Módulos cocientes $\mathcal G$ de $\mathcal F'$, planos sobre
$S'$, cuyo soporte está contenido en $U'$. Manifiestamente, $A'$ está
representado por un subconjunto abierto del preesquema
$\operatorname{Quot}_{\mathcal F/X/S}$ que representa $A$. De ello resulta
que los Teoremas (3.1) y (3.2)
\src{267}
siguen siendo válidos cuando se supone que $X$ es cuasiproyectivo sobre $S$ en
lugar de proyectivo sobre $S$, siempre que en las conclusiones también se
sustituya «proyectivo» por «cuasiproyectivo» y que
$\underline{\operatorname{Quot}}_{\mathcal F/X/S}(S')$ denote ahora el
conjunto de cocientes coherentes $\mathcal G$ de $\mathcal F'$, planos sobre
$S'$, cuyo soporte es propio sobre $S'$.

\medskip
\noindent b. Más generalmente, pueden imponerse a los cocientes $\mathcal G$
de $\mathcal F'=\mathcal F\otimes_{\mathcal O_S}\mathcal O_{S'}$, planos
sobre $S'$, toda clase de condiciones naturales adicionales estables por
cambio de base, obteniéndose así subfuntores de
$\underline{\operatorname{Quot}}_{\mathcal F/X/S}$ cuya representabilidad
puede buscarse. En muchos casos, el criterio usual demuestra que estos
funtores vuelven a estar representados por subconjuntos abiertos de
$\operatorname{Quot}_{\mathcal F/X/S}$. Esto ocurre, en particular, cuando se
impone una de las propiedades adicionales siguientes:

\smallskip
\noindent $1^\circ$ Las dimensiones de los ciclos primos asociados a los
Módulos $\mathcal G_{s'}$ ($s'\in S'$) inducidos sobre las fibras $X'_{s'}$
pertenecen a un conjunto dado de enteros.

\smallskip
\noindent $2^\circ$ (Cuando $\mathcal F=\mathcal O_X$, de modo que
$\mathcal G$ corresponde a un subpreesquema cerrado $Y$ de $X'$): $Y$ es un
preesquema liso ([3], IV) sobre $S'$, respectivamente normal sobre $S'$ (es
decir, las fibras $Y_{s'}$ son normales «sobre $k(s')$», lo que significa
normales después de toda extensión del cuerpo base), respectivamente (cuando
$X$ es plano sobre $S$) una intersección completa local relativa en $X'$ sobre
$S'$ (es decir, las fibras $Y_{s'}$ son intersecciones completas locales en los
$X'_{s'}$).\footnote{El texto impreso omite las primas en estas propiedades de $Y$ sobre la base variable $S'$ y dentro de $X'$. La hipótesis relativa al $X$ original plano sobre $S$ permanece sin primas.}

Otras condiciones harían intervenir propiedades cohomológicas de los Módulos
$\mathcal G_{s'}$ inducidos sobre los $X'_{s'}$, etc. Por supuesto, una
conjunción de condiciones, cada una de ellas representada por un subconjunto
abierto $U_i$ de $\operatorname{Quot}_{\mathcal F/X/S}$, está representada por
la intersección de estos subconjuntos abiertos. Por ejemplo, asignar a todo
$S'$ sobre $S$ el conjunto de subesquemas cerrados $Y$ de
$X'=X\times_SS'$ que son recubrimientos étale ([3], I) de $S'$ de rango dado
$r$ da un funtor contravariante representable de $S'$.

\medskip
\noindent c. Los preesquemas $\operatorname{Hom}_S(X,Y)$,
$\Pi_{X/S}(Z/X)$ y $\operatorname{Isom}_S(X,Y)$ definidos en [2], II, C,
n.º 2 existen bajo hipótesis de proyectividad convenientes y aparecen como
subconjuntos abiertos de preesquemas de Hilbert adecuados. Puesto que
\[
\operatorname{Hom}_S(X,Y)=\Pi_{X/S}((X\times_SY)/X),
\]
el caso de $\operatorname{Hom}_S(X,Y)$ se reduce al de
$\Pi_{X/S}(Z/X)$. Para todo $S'$ sobre $S$, el conjunto de secciones de
$Z'=Z\times_SS'$ sobre $X'=X\times_SS'$ está en correspondencia biyectiva con
el conjunto de subpreesquemas $\Gamma$ de $Z'$ (necesariamente cerrados si $Z$
es separado sobre $X$) tales que el morfismo $\Gamma\longrightarrow X'$
inducido por $Z'\longrightarrow X'$ sea un isomorfismo. De este modo,
\src{268}
cuando $X$ es plano y propio sobre $S$ y $Z$ es cuasiproyectivo sobre $S$,
$\Pi_{X/S}(Z/X)$ existe y aparece como subpreesquema abierto de
$\operatorname{Hilb}_{Z/S}$. Por consiguiente, cuando $X$ es proyectivo y
plano sobre $S$ e $Y$ es cuasiproyectivo sobre $S$,
$\operatorname{Hom}_S(X,Y)$ existe y aparece como subpreesquema abierto de
$\operatorname{Hilb}_{(X\times_SY)/S}$. Cuando $X$ e $Y$ son ambos
proyectivos sobre $S$, se sigue inmediatamente que
$\operatorname{Isom}_S(X,Y)$ también existe y está representado por un
subconjunto abierto de $\operatorname{Hom}_S(X,Y)$. Análogamente, cuando $X$
es plano y proyectivo sobre $S$ e $Y$ es cuasiproyectivo sobre $S$, el
$S$-preesquema $\operatorname{Imm}_S(X/Y)$ correspondiente al subfuntor de
$\operatorname{Hom}_S(X,Y)$ cuyos elementos son los $S'$-homomorfismos
$X'\longrightarrow Y'$ que son inmersiones también está representado por un
subconjunto abierto de $\operatorname{Hom}_S(X,Y)$.

Sea $\mathcal L$ (respectivamente, $\mathcal M$) un haz invertible sobre $X$
(respectivamente, $Y$), muy amplio relativamente a $S$. Entonces
$\mathcal L\otimes_{\mathcal O_S}\mathcal M$ es un haz muy amplio sobre
$X\times_SY$ relativamente a $S$. Por consiguiente, para todo polinomio $P$
con coeficientes racionales, $\operatorname{Hilb}^P_{(X\times_SY)/S}$ está
definido y es un preesquema cuasiproyectivo sobre $S$. Induce, por tanto,
sobre $\operatorname{Hom}_S(X,Y)$ un subconjunto abierto y cerrado,
cuasiproyectivo sobre $S$, que denotaremos por
$\operatorname{Hom}_S(X,Y)^P$. Así, las secciones de
$\operatorname{Hom}_S(X,Y)^P$ sobre $S$ son los $S$-morfismos
$g:X\longrightarrow Y$ tales que, para todo entero $n$,
\[
\chi\!\left(\left(\mathcal L\otimes_{\mathcal O_X}g^*(\mathcal M)\right)^{
\otimes n}\right)=P(n).
\]

De este modo se obtienen generalizaciones del teorema de Matsusaka que afirma
que los automorfismos de una variedad proyectiva «polarizada» forman un grupo
algebraico. Aquí la afirmación tiene un sentido manifiestamente más preciso,
pues este grupo se define como solución de un problema universal. Obsérvese
también que, sobre un cuerpo algebraicamente cerrado, el grupo de
automorfismos considerado anteriormente se obtiene del «verdadero» grupo
definido aquí pasando al cociente por los elementos nilpotentes. Esto explica
por qué las construcciones antiguas tienen pocas posibilidades de funcionar
sobre un cuerpo base no perfecto: el ideal de elementos nilpotentes que
aparece después de una extensión del cuerpo base no tiene por qué estar
«definido sobre $k$». La misma observación se aplica a la mayoría de las
construcciones de estilo antiguo.

\fsection{5}{Estudio infinitesimal de los esquemas de Hilbert}

Se deduce del resultado siguiente:

\src{269}
\result{PROPOSICIÓN}{5.1} Sea $S$ un preesquema, sea $S_0$ un subpreesquema
definido por un ideal cuasicoherente $\mathcal J$ de cuadrado nulo, sea $X$ un
$S$-preesquema y sea $\mathcal F$ un Módulo cuasicoherente sobre $X$. Pongamos
$X_0=X\times_SS_0$ y
$\mathcal F_0=\mathcal F\otimes_{\mathcal O_S}\mathcal O_{S_0}$, y sea
$\mathcal G_0=\mathcal F_0/\mathcal K_0$ un Módulo cuasicoherente cociente de
$\mathcal F_0$, plano sobre $S_0$. Para todo subconjunto abierto $U$ de $X$,
sea $\mathcal E(U)$ el conjunto de Módulos cuasicoherentes cocientes
$\mathcal G$ de $\mathcal F|U$, planos sobre $S$, tales que
$\mathcal G\otimes_{\mathcal O_S}\mathcal O_{S_0}=\mathcal G_0$. Al variar
$U$, los $\mathcal E(U)$ son, pues, las secciones de un haz $\mathcal E$ sobre
$X$.\footnote{El texto impreso lleva $U$; al variar los abiertos $U$ de $X$, los conjuntos $\mathcal E(U)$ definen un haz sobre $X$.} El haz de grupos
\[
\mathcal A=\operatorname{Hom}_{\mathcal O_{X_0}}
(\mathcal K_0,\mathcal G_0\otimes_{\mathcal O_{S_0}}\mathcal J)
\]
actúa naturalmente sobre $\mathcal E$, que se convierte así en un haz
«formalmente homogéneo principal bajo $\mathcal A$» (es decir, para todo
abierto $U$ de $X$, $\mathcal E(U)$ es vacío o es un conjunto homogéneo
principal bajo $\mathcal A(U)$).

De ello se sigue que:

\result{COROLARIO}{5.2} Supongamos que localmente sobre $X$ existe una
extensión $\mathcal G$ de $\mathcal G_0$ a un cociente de $\mathcal F$ plano
sobre $S$ (es decir, que las fibras del haz $\mathcal E$ no son vacías).
Existe entonces una clase de obstrucción canónica
\[
c(\mathcal G_0)\in H^1(X,\mathcal A)
\]
cuya anulación es necesaria y suficiente para que exista una extensión global
$\mathcal G$ de $\mathcal G_0$ a un cociente de $\mathcal F$ plano sobre $S$.
Si esta clase se anula, el conjunto $\mathcal E(X)$ de todas las extensiones
posibles es un conjunto homogéneo principal bajo
\[
\mathcal A(X)=\operatorname{Hom}_{\mathcal O_{X_0}}
(\mathcal K_0,\mathcal G_0\otimes_{\mathcal O_{S_0}}\mathcal J).
\]
\footnote{En esta última fórmula, el texto impreso lleva $\mathcal O_X$ en
lugar de $\mathcal O_{X_0}$, aunque la definición inmediatamente anterior de
$\mathcal A$ fija $\mathcal O_{X_0}$.}
En particular, existe una extensión global si $H^1(X,\mathcal A)=0$.

\result{COROLARIO}{5.3}\footnote{El texto impreso numera este corolario como
«3.3»; sigue a (5.1) y (5.2) en el Apartado 5.} Supongamos que
$Q=\operatorname{Quot}_{\mathcal F/X/S}$ existe (cf. 4 (a)); por ejemplo, que
$X$ es cuasiproyectivo sobre un $S$ localmente noetheriano y que $\mathcal F$
es coherente. Sea $x\in Q$ correspondiente a una extensión residual
$K=k(x)$ de $k(s)$ ($s\in S$):
\src{270}
así, $x$ está definido por un Módulo coherente cociente
$\mathcal G_0=\mathcal F_0/\mathcal K_0$ del Módulo
$\mathcal F_0=\mathcal F_K$ sobre el $K$-preesquema $X_K$. Sea $\mathcal A$
el haz coherente sobre $X_K$ definido por
\[
\mathcal A=\operatorname{Hom}_{\mathcal O_{X_0}}(\mathcal K_0,\mathcal G_0).
\]
Entonces el espacio tangente de Zariski a la fibra $Q_s$ en $x$ (el dual
sobre $K$ de $\mathfrak m/\mathfrak m^2$, donde $\mathfrak m$ es el ideal
maximal de $\mathcal O_{Q_K,x}$) es canónicamente isomorfo a
$H^0(X_K,\mathcal A)$.

La descripción del espacio tangente de Zariski admite la generalización
siguiente. Dado un $S$-morfismo $g:S'\longrightarrow Q$, es decir, una sección
$g'$ de $Q'=Q\times_SS'$ sobre $S'$, el Módulo
\[
\Omega=g^*(\Omega^1_{Q/S})=g'^*(\mathcal J/\mathcal J^2)
\]
(donde $\mathcal J$ es el ideal sobre $Q'$ definido por la sección $g'$ de
$Q'$ sobre $S'$) está caracterizado por la fórmula siguiente, funtorial en el
Módulo coherente $\mathcal M$ sobre $S'$:
\[
\operatorname{Hom}_{\mathcal O_{S'}}(\Omega,\mathcal M)
\simeq H^0(X',\mathcal A),
\]
donde $\mathcal A$ es el Módulo sobre $X'=X\times_SS'$ definido por
\[
\mathcal A=\operatorname{Hom}_{\mathcal O_{X'}}
(\mathcal K,\mathcal G\otimes_{\mathcal O_{S'}}\mathcal M).
\]
\footnote{El texto impreso pone $\mathcal O_S$ en este producto tensorial;
aquí $\mathcal G$ y $\mathcal M$ son Módulos sobre $S'$, y la construcción de
deformación de cuadrado nulo exige el producto tensorial sobre
$\mathcal O_{S'}$.}
Aquí $\mathcal G=\mathcal F'/\mathcal K$ es el Módulo cociente de
$\mathcal F'=\mathcal F\otimes_{\mathcal O_S}\mathcal O_{S'}$ que corresponde
a $g$. En efecto, aplíquese (5.1) sustituyendo $S_0$ por $S'$ y $S$ por el
preesquema $D(\mathcal M)=(S',\mathcal O_{S'}+\mathcal M)$, donde
$\mathcal M$ se considera como un ideal de cuadrado nulo.

Cuando $\mathcal F=\mathcal O_X$ en (5.1), el dato de $\mathcal G_0$
corresponde a un subpreesquema cerrado $Y_0$ de $X_0$, plano sobre $S_0$,
definido por el ideal $\mathcal J_0=\mathcal K_0$. La fórmula (*) da entonces
\[
\mathcal A=\operatorname{Hom}_{\mathcal O_{X_0}}
(\mathcal J_0/\mathcal J_0^2,
\mathcal O_{Y_0}\otimes_{\mathcal O_{S_0}}\mathcal J),
\]
donde $\mathcal J_0/\mathcal J_0^2$\footnote{El texto impreso omite el subíndice $0$ en las dos menciones del haz conormal y en $\mathcal J_0^{-1}$ más abajo. El Ideal de $Y_0$ en $X_0$ es $\mathcal J_0$; el Ideal $\mathcal J$ de la base es distinto.} se interpreta como el haz conormal de $Y_0$
en $X_0$, denotado también por
\src{271}
$\check{\mathcal N}_{Y_0/X_0}$. Conviene entonces considerar $\mathcal A$
como un Módulo sobre $Y_0$ y calcular $H^0$ y $H^1$ sobre $Y_0$. Además, si
$Y_0$ es localmente una intersección completa en $X_0$, con $X$ plano sobre
$S$, queda garantizada la existencia de extensiones locales en (5.1); asimismo,
$\mathcal J_0/\mathcal J_0^2$ es localmente libre sobre $Y_0$, y puede escribirse
\[
\mathcal A=\mathcal N_{Y_0/X_0}\otimes_{\mathcal O_{S_0}}\mathcal J,
\]
donde el primer factor de la derecha es el haz normal de $Y_0$ en $X_0$.
Utilizando el criterio fundamental de lisura ([3], III, 3.1), se obtiene, por
ejemplo:

\result{COROLARIO}{5.4} Bajo las hipótesis de (5.3), supongamos que
$\mathcal F=\mathcal O_X$, que $X$ es plano sobre $S$ y que el subpreesquema
cerrado $Y_0$ de $X_0$ correspondiente a $\mathcal G_0$ es localmente una
intersección completa. Entonces el espacio tangente de Zariski a $Q_s$ en $x$
es canónicamente isomorfo a $H^0(Y_0,\mathcal N_{Y_0/X_0})$. Si
$H^1(Y_0,\mathcal N_{Y_0/X_0})=0$, entonces el preesquema de Hilbert
$\operatorname{Hilb}_{X/S}$ es liso sobre $S$ en $x$.
($\mathcal N_{Y_0/X_0}$ es el haz normal de $Y_0$ en $X_0$.)

\result{OBSERVACIÓN}{5.5} Este enunciado se aplica, en particular, cuando
$Y_0$ es una intersección completa en $X_0$ definida por una ecuación, es
decir, un «divisor de Cartier» positivo. Entonces $\mathcal N_{Y_0/X_0}$ es
isomorfo al haz sobre $Y_0$ inducido por el haz invertible
$\mathcal J_0^{-1}$ sobre $X_0$ definido por el divisor $Y_0$. Esta situación
se presenta especialmente en el estudio de familias de divisores positivos
sobre una variedad proyectiva no singular $X_0$. El isomorfismo entre el
espacio tangente de Zariski en el punto $x$ de $Q$ (o, equivalentemente, del
subconjunto abierto $D$ de $Q$ correspondiente a los divisores) y
$H^0(Y_0,\mathcal N_{Y_0/X_0})$ era conocido en geometría algebraica clásica
como el «homomorfismo característico» del primero al segundo. Solo estaba
definido cuando $x$ era un punto liso de la variedad de parámetros $T$ de una
«familia continua completa» de divisores, es decir, desde nuestro punto de
vista, de una componente irreducible del esquema $D$ provisto de su estructura
reducida inducida. El
espacio tangente a $T$ en $x$ es entonces un subespacio del espacio tangente a
$D$ en $x$, por lo que el homomorfismo característico clásico es, en efecto,
inyectivo, pero solo es sobreyectivo bajo condiciones adicionales; por ejemplo,
cuando $D$ es íntegro en $x$. De hecho, Zappa [8] construyó un ejemplo
\src{272}
(con $X$ una superficie proyectiva no singular sobre los números complejos) en
el que el homomorfismo característico no es sobreyectivo ni siquiera en el
punto genérico de $T$. Así, $D$ no es íntegro ni siquiera en el punto genérico
de la componente irreducible considerada. Esto demuestra de manera
especialmente llamativa hasta qué punto las variedades con elementos
nilpotentes son necesarias para comprender fenómenos de la teoría más clásica
de superficies.\footnote{La fe de erratas de 1962 añade que Mumford acababa de
construir una componente irreducible del esquema de Hilbert de
$\mathbf P^3_{\mathbf C}$ cuyos puntos generales representan curvas no
singulares de grado 14 y género 24, que es no reducida en su punto genérico;
al hacer explotar una de estas curvas se obtiene también una variedad triple
proyectiva regular cuyo esquema formal de módulos es genéricamente no
reducido.}

\subsection*{5.6}

En (5.4) dimos un criterio de lisura aplicable, en particular, a esquemas de
divisores. Kodaira dio en [6] un criterio diferente, a saber, la anulación de
$H^1(X_0,\mathcal L)$, donde $\mathcal L=\mathcal J_0^{-1}$ es el haz
invertible sobre $X_0$ definido por el divisor $Y_0$. Este criterio es válido
cuando $S$ es el espectro de un cuerpo de característica cero, y fue demostrado
en [6] por métodos trascendentes cuando el cuerpo base es $\mathbf C$. Más
generalmente, siendo de nuevo $S$ arbitrario, la condición de Kodaira basta
para que el morfismo canónico
\[
D\longrightarrow\operatorname{Pic}_{X/S}
\]
del preesquema de divisores al preesquema de Picard de $X/S$ sea liso en el
punto $x$ considerado (como se deduce fácilmente del criterio usual de lisura
una vez conocida la existencia de $\operatorname{Pic}_{X/S}$). Si, además,
$\operatorname{Pic}_{X/S}$ es liso sobre $S$ en la imagen de $x$ (por ejemplo,
si $\operatorname{Pic}_{X/S}$ es liso sobre $S$), entonces $D$ es liso sobre
$S$ en $x$. Por otra parte, Cartier demostró que todo preesquema en grupos
localmente de tipo finito sobre un cuerpo $k$ de característica cero es liso
sobre $k$. Combinando estos dos resultados se recupera el resultado de
Kodaira. Estas observaciones demuestran también que, sobre un cuerpo $K$ de
característica $p>0$, cuando $\operatorname{Pic}_{X/S}$ no es liso sobre $K$\footnote{El texto impreso lleva $k$; el cuerpo base de esta frase en característica positiva es $K$.}
(como cuando $X$ es la superficie de Igusa), la condición
$H^1(X_0,\mathcal L)=0$ implica, por el contrario, que $D$ no es liso en $x$,
y que incluso es no reducido en $x$ si $K$ es algebraicamente cerrado.

Finalmente, registremos el resultado siguiente, que desempeña un papel
importante en el estudio infinitesimal de los espacios fibrados:

\result{PROPOSICIÓN}{5.7} Sea $X$ un preesquema finito y plano sobre un $S$
localmente noetheriano, y sea $Z$ un preesquema sobre $X$ tal que
$\Pi_{X/S}(Z/X)$ exista (como ocurre si $Z$ es cuasiproyectivo sobre $X$). Si
$Z$ es liso sobre $X$, entonces $\Pi_{X/S}(Z/X)$ es liso sobre $S$.

Esto resulta inmediatamente de la definición y del criterio usual de lisura
([3], III, 3.1). Obsérvese que, cuando $X$ es finito y plano sobre $S$, la
existencia de $\Pi_{X/S}(Z/X)$\footnote{El texto impreso lleva aquí $Z/K$; no
se ha definido ningún $K$, y el funtor de secciones considerado desde (5.7) es
$\Pi_{X/S}(Z/X)$.} puede tratarse de manera muy elemental, sin
\src{273}
utilizar la teoría de los esquemas de Hilbert. Por ejemplo, si $X$ es radicial
sobre $S$, entonces $\Pi_{X/S}(Z/X)$ existe sin ninguna restricción sobre $Z$.
Como otro ejemplo, sea $T$ un $S$-preesquema y sea $T_n$ el «entorno
infinitesimal de orden $n$» de la diagonal de $T\times_ST$ en $T\times_ST$,
provisto de los morfismos $p_1,p_2:T_n\longrightarrow T$ inducidos por las dos
proyecciones. Consideremos $T_n$ como un preesquema finito sobre $T$ por medio
de $p_1$, y supongamos además que es plano sobre $T$ (como ocurre cuando $T$
es liso sobre $S$). Para todo preesquema $X$ sobre $T$, pongamos
\[
(X/T/S)^{(n)}=
\Pi_{T_n/T}\bigl(p_2^*(X/T)/T_n\bigr).
\]
\footnote{El texto impreso coloca $T_n/S$ bajo el símbolo $\Pi$. La
Proposición 5.7 se aplica aquí al morfismo finito y plano
$p_1:T_n\to T$, y el resultado es explícitamente un preesquema sobre $T$; por
tanto, el denominador es $T_n/T$.}
Este es un preesquema sobre $T$ llamado \emph{fibrado de gérmenes de secciones
de orden $n$ de $X$ sobre $T$} (relativo a $S$). Depende funtorialmente de $X$
y es liso sobre $T$ si $X$ lo es.

\fsection{6}{Relaciones con la norma y los productos simétricos}

Sea $S$ un preesquema, sean $X$ e $Y$ $S$-preesquemas, y sea
\[
u:(X/S)^n\longrightarrow Y
\]
un $S$-morfismo simétrico de la $n$-ésima potencia cartesiana de $X/S$ en
$Y$. Para simplificar, supongamos que $S$ es localmente noetheriano y que $X$
e $Y$ son de tipo finito sobre $S$. A todo Módulo coherente $\mathcal F$ sobre
$X$ cuyo soporte es finito sobre $S$, que es plano sobre $S$ y tiene rango
relativo $n$ sobre $S$ (es decir, $f_*(\mathcal F)$ es un Módulo localmente
libre de rango $n$ sobre $S$), puede asociársele entonces de manera natural una
sección de $Y$ sobre $S$:
\[
\pi^u_{X/S}(\mathcal F)\in\Gamma(Y/S).
\]
No daremos aquí la definición formal, sino que nos limitaremos a observar que
el formalismo resultante es una generalización natural del formalismo usual de
normas y trazas. Para todo $X/S$ separado, la ley de norma tiene como blanco canónico el espacio de potencias divididas $\Gamma^n_S(X)$. Si $X$ es plano sobre $S$ y existe la $n$-ésima potencia simétrica (por ejemplo, cuando las órbitas de $\mathfrak S_n$ en $(X/S)^n$ están contenidas en abiertos afines), el morfismo canónico $\operatorname{Symm}^n_S(X)\to\Gamma^n_S(X)$ es un isomorfismo, y se obtiene un elemento canónico
\[
\pi_{X/S}(\mathcal F)\in\Gamma(\operatorname{Symm}^n_S(X)/S),
\]
\src{274}
del cual pueden recuperarse los $\pi^u_{X/S}(\mathcal F)$. Otro caso
importante es aquel en que $X$ es un monoide conmutativo sobre $S$, $X=Y$, y
$u$ procede de la ley de composición de $X$. Escribimos entonces simplemente
$\pi(\mathcal F)$ para la sección de $X$ sobre $S$ asociada al Módulo
$\mathcal F$ sobre $X$.

Supongamos ahora que $\mathcal F$ es un Módulo coherente sobre $X$ tal que
$\operatorname{Quot}_{\mathcal F/X/S}$ existe, o al menos tal que el funtor
$\underline{\operatorname{Quot}}^n_{\mathcal F/X/S}$ que asigna a todo $S'$
sobre $S$ el conjunto de haces cocientes coherentes $\mathcal M$ de
$\mathcal F'=\mathcal F\otimes_{\mathcal O_S}\mathcal O_{S'}$ que son planos
sobre $S'$ y tienen rango relativo $n$ está representado por un $S$-preesquema
$\operatorname{Quot}^n_{\mathcal F/X/S}$. (Cuando $X$ es cuasiproyectivo
sobre $S$, $\operatorname{Quot}^n_{\mathcal F/X/S}$ existe y, con las
notaciones del Apartado 3, es $\operatorname{Quot}^P_{\mathcal F/X/S}$, donde
$P$ es el polinomio constante $n$.) Como la formación de
$\pi^u_{X/S}(\mathcal M)$ es compatible con el cambio de base, se obtiene un
morfismo canónico
\[
\pi^u_{X/S}:\operatorname{Quot}^n_{\mathcal F/X/S}\longrightarrow Y,
\]
y, en general,
\[
\pi_{X/S}:\operatorname{Quot}^n_{\mathcal F/X/S}
\longrightarrow\Gamma^n_S(X).
\]
Cuando $X$ es plano sobre $S$, identificamos este blanco con $\operatorname{Symm}^n_S(X)$.
Independientemente de la identificación precedente, cuando $X$ es separado sobre $S$, la construcción mediante la familia finita y plana para $\mathcal F=\mathcal O_X$ da un morfismo
\[
\pi_{X/S}:\operatorname{Hilb}^n_{X/S}
\longrightarrow\operatorname{Symm}^n_S(X).
\]
Este es manifiestamente un isomorfismo para $n=0$ o $n=1$. Pero para
$n>1$\footnote{El texto impreso repite $n\geq1$, contradiciendo el
isomorfismo que acaba de afirmarse para $n=1$.}, incluso cuando $S$ es el
espectro de un cuerpo $k$ y $X$ es liso sobre $S$, no es en general un
isomorfismo, ni siquiera un morfismo inyectivo: un subesquema de dimensión cero
de $X$ (correspondiente, por ejemplo, a un ideal $\mathfrak i$ primario para el
ideal maximal en un anillo local $\mathcal O_{X,x}$, para un punto cerrado $x$
de $X$) no está determinado por el ciclo que define (en este caso, por la
codimensión sobre $k$ de $\mathfrak i$ en $\mathcal O_{X,x}$). Solo puede
afirmarse lo siguiente, siendo de nuevo $S$ arbitrario:

\medskip
\noindent a. Si $X$ es liso sobre $S$, el morfismo norma induce un isomorfismo
del subconjunto abierto de $\operatorname{Hilb}^n_{X/S}$ que clasifica los
recubrimientos étale de rango $n$ contenidos en $X$ (cf. n.º 4, (b)) sobre el
subconjunto abierto de $\operatorname{Symm}^n_S(X)$ correspondiente a los
$n$-ciclos sin componentes múltiples.
\src{275}

\medskip
\noindent b. Si, además, $X$ tiene dimensión relativa 1 sobre $S$, entonces el
morfismo norma es un isomorfismo de $\operatorname{Hilb}^n_{X/S}$ sobre
$\operatorname{Symm}^n_{X/S}$.

Esto se debe a que un subesquema de dimensión cero de una curva algebraica no
singular está determinado por su ciclo correspondiente. La misma observación
se aplica, más generalmente, a los divisores de Cartier positivos sobre un
esquema algebraico no singular (y es posible que, en este caso tan especial,
la variedad de Chow dé el mismo objeto que el esquema de Hilbert).

\fsection{7}{Complementos y cuestiones}

Como observó J.-P. Serre, un ejemplo bien conocido de Nagata demuestra que
pueden encontrarse un esquema $S$, espectro de un cuerpo $k$; un $S$-esquema
$S'$, espectro de una extensión cuadrática $k'$ de $k$; y un $S'$-esquema
propio y liso (pero no proyectivo) $X$ de dimensión 3 tal que
$\Pi_{S'/S}(X/S')$ no existe. Se sigue, a fortiori, que el esquema de
Hilbert $\operatorname{Hilb}^2_{X/S}$ no existe (ni siquiera el $k$-esquema
que representaría los recubrimientos étale de $S$ de rango 2 contenidos en
$X$, y a fortiori tampoco el cuadrado simétrico de $X$; cf. el apartado
anterior). Esto impone serias limitaciones a la posibilidad de realizar
construcciones no proyectivas en geometría algebraica. (Sin embargo, es
plausible que tales limitaciones no aparezcan en geometría analítica, como
tampoco aparecen en geometría formal (cf. [2], II).) Por otra parte, si $X$ es
un esquema propio sobre $S=\operatorname{Spec}k$ y $Z$ es cuasiproyectivo
sobre $X$, entonces $\Pi_{X/S}(Z/X)$ existe y es un esquema que es unión
disjunta de una sucesión de esquemas cuasiproyectivos sobre $S$ (como en el
caso proyectivo (3.1)). En efecto, se reduce al caso en que el propio $X$ es
proyectivo dominando $X$ por un $S$-esquema proyectivo $X'$. No daremos los
detalles de la demostración, que utiliza también el resultado de factorización
para un morfismo finito recogido en [2], I, A, n.º 2 (b). El método funciona
porque, siendo $S$ el espectro de un cuerpo, el $X'$ que interviene en el lema
de Chow es automáticamente plano sobre $S$.

No sé si el resultado sigue siendo cierto sin hipótesis sobre $S$, suponiendo
únicamente que $X$ es propio y plano sobre $S$ y que $Z$ es cuasiproyectivo
sobre $X$. Un caso importante en las aplicaciones es aquel en que $Z$ es un
subesquema cerrado de $X$; si $\Pi_{X/S}(Z/X)$ existe, es necesariamente un
subesquema cerrado de $S$. Puede
\src{276}
construirse directa y muy sencillamente cuando $X$ es proyectivo sobre $S$,
sin utilizar la teoría de los esquemas de Hilbert. Más generalmente, el mismo
método demuestra que, si $Z$ es afín sobre $X$, entonces $\Pi_{X/S}(Z/X)$
existe y es afín sobre $S$. También demuestra que, si $X$ es propio y plano
sobre $S$ (no necesariamente proyectivo sobre $S$), entonces, para todo
fibrado vectorial localmente trivial $Z$ sobre $X$, $\Pi_{X/S}(Z/X)$ existe y
es un fibrado vectorial sobre $S$. Sería deseable retomar y unificar estos
resultados.

\section*{BIBLIOGRAFÍA}
\addcontentsline{toc}{section}{Bibliografía}

\begin{description}
\item[{[1]}] BOREL (A.) y SERRE (J.-P.). -- Le théorème de Riemann--Roch,
\emph{Bull. Soc. math. France}, vol. 86, 1958, pp. 97--136.

\item[{[2]}] GROTHENDIECK (Alexander). -- Technique de construction et
théorèmes d'existence en géométrie algébrique, I, II, III, \emph{Séminaire
Bourbaki}, vol. 13, 1959/60, n.º 190, 29 pp., n.º 195, 22 pp.; y vol. 14,
1960/61, n.º 212, 20 pp.

\item[{[3]}] GROTHENDIECK (Alexander). -- Séminaire de Géométrie algébrique,
I, II, III, IV. -- París, Institut des hautes Études scientifiques, 1960/61
(mimeografiado).

\item[{[4]}] GROTHENDIECK (Alexander). -- Technique de construction en
géométrie analytique, IV, V, \emph{Séminaire Cartan}, vol. 13, 1960/61, n.os 11
y 12.

\item[{[5]}] GROTHENDIECK (A.) y DIEUDONNÉ (J.). -- \emph{Éléments de
géométrie algébrique}, I: Le langage des schémas. -- París, Presses
universitaires de France, 1960 (Institut des hautes Études scientifiques,
Publications mathématiques, 4); II, III, IV (de próxima aparición).

\item[{[6]}] KODAIRA (K.). -- Characteristic linear systems of complete
continuous systems, \emph{Amer. J. of Math.}, vol. 78, 1956, pp. 716--744.

\item[{[7]}] SERRE (Jean-Pierre). -- Exemples de variétés projectives en
caractéristique $p$ non relevables en caractéristique 0, \emph{Proc. Nat. Acad.
Sc. U.S.A.}, vol. 47, 1961, pp. 108--109.

\item[{[8]}] ZAPPA (G.). -- Sull'esistenza, sopra le superficie algebriche, di
sistemi continui completi infiniti, la cui curva è a serie caratteristica
incompleta, \emph{Pont. Acad. Sc. Acta}, vol. 9, 1945, pp. 91--93.
\end{description}

\section*{ADITIVO}
\addcontentsline{toc}{section}{Aditivo}
\begin{center}
{\small [añadido durante la corrección de pruebas]}
\end{center}

Parece ahora que las conjeturas de [2], III, n.º 8 son falsas, incluso para
variedades no singulares sobre un cuerpo de característica cero, tanto en lo
que respecta a la existencia como a la cuasiproyectividad del cociente, e
incluso cuando $G$ actúa con grafo cerrado.
